\documentclass[11pt]{article}
\usepackage[margin=0.7in]{geometry}
\usepackage{titlesec}
\usepackage{enumitem}
\usepackage{parskip}

\titleformat{\section}{\Large\bfseries}{\thesection}{1em}{}
\titleformat{\subsection}{\large\bfseries}{\thesubsection}{1em}{}

\setlength{\parskip}{0.3em}

\title{\textbf{Week 5: Hard Bop (Mid-1950s to Early 1960s)} \\ \large Quick Reference}
\author{}
\date{}

\begin{document}

\maketitle

\section*{Period Overview}

\textbf{What is Hard Bop?} Bebop + soul/gospel/blues. Mid-1950s reaction against cool jazz's "whitening." East coast sound combining bebop's technical complexity with accessible grooves, gospel influences, and blues feeling. Made jazz danceable again to compete with rock and roll. Key figures: Art Blakey, Horace Silver, Sonny Rollins, Cannonball Adderley.

\textbf{Key Sound:} Funky, bluesy, groove-oriented, soulful. Strong driving groove, walking bass, powerful interactive drumming.

\textbf{Relation to Bebop:} Maintains bebop's complex chord changes and virtuosity but adds soul, groove, and accessibility. Walked harmonic complexity down slightly in favor of feeling.

\section*{Required Listening (8 Pieces)}

\subsection*{Art Blakey - "Blues March" (1958)}
This hard bop piece exemplifies the mid-1950s East Coast reaction against cool jazz by combining bebop virtuosity with accessible gospel-influenced grooves. The instrumentation features Lee Morgan's bright trumpet, Benny Golson's warm tenor saxophone, and Blakey's thunderous interactive drumming with his signature press roll between solos. The medium tempo 12-bar blues form uses a march rhythm that emphasizes beats 1 and 3 rather than jazz's typical 2 and 4, creating a more danceable feel. While maintaining bebop's complex chord changes and technical virtuosity during solos, the piece adds gospel call-and-response patterns and a strong walking bass groove. This represents hard bop's mission to make jazz accessible to wider audiences while keeping bebop's sophistication, competing with the rise of rock and roll.

\subsection*{Horace Silver - "Song for My Father" (1964)}
This soul jazz piece exemplifies the funky gospel-influenced substyle of hard bop through its Cape Verdean/Caribbean musical influences and accessible vamp structure. The instrumentation features Silver's percussive funky piano comping style over a Latin/bossa nova beat, while the bass alternates just two notes rather than walking. The simple catchy melody and repetitive vamp structure simplify bebop's complex harmony in favor of groove and positive celebratory energy. Silver's compositional approach reflects the "subconsciously racial" quality of hard bop discussed in the readings, drawing from Black folk, church, and Caribbean traditions. This piece demonstrates how soul jazz prioritized feeling and accessibility over bebop's intellectual complexity while maintaining jazz sophistication.

\subsection*{Sonny Rollins - "St. Thomas" (1956)}
This hard bop piece features a calypso rhythm and Caribbean folk melody that Rollins learned from his mother, demonstrating personal cultural expression in jazz improvisation. Rollins' warm "thick" tenor tone shows Coleman Hawkins' influence while his economical fragmented solos reflect Count Basie's approach to space and phrasing. The instrumentation features bass alternating notes rather than strict walking bass and drums tuned melodically to complement the Caribbean feel at medium tempo. Rollins transfers Charlie Parker's bebop fluidity and virtuosity to tenor saxophone but replaces bebop's straight-ahead swing with a relaxed accessible calypso groove. This exemplifies hard bop's balance between maintaining bebop's technical complexity and adding cultural influences and accessibility for broader audiences.

\subsection*{Miles Davis - "Boplicity" from \emph{Birth of the Cool} (1949)}
This cool jazz piece represents a direct reaction against bebop's intensity and frenetic energy through its restrained orchestral approach. The 9-piece nonet instrumentation includes French horn and tuba, creating lush orchestral brass textures unlike bebop's typical small group format. Miles plays mellow muted trumpet at a slower tempo with gentle restrained dynamics, "shaking ears more gently" compared to Parker's "furious" bebop playing. Gil Evans' sophisticated arrangement integrates improvisation within composed passages, emphasizing melodic beauty over virtuosic complexity. Hard bop later emerged as a reaction against this cool jazz "whitening" of bebop, reclaiming the music's Black cultural roots through soul and gospel influences.

\subsection*{Miles Davis - "Oleo" from \emph{Relaxin'} (1956)}
This pure hard bop piece is a contrafact using fast "I Got Rhythm" chord changes, demonstrating bebop's harmonic vocabulary with hard bop's new sensibility. The instrumentation features Miles' distinctive Harmon mute, John Coltrane's maturing tenor sound, and Paul Chambers' standout walking bass at breakneck tempo. Recorded during the famous Prestige sessions where Miles cut four albums in two days with mostly first takes, the performance shows spontaneous interactive improvisation. The piece maintains bebop's complex chord progressions and virtuosic soloing but adds Miles' cool restraint and more accessible phrasing. This represents the transitional moment where bebop's technical demands merged with hard bop's emphasis on groove and feeling.

\subsection*{Cannonball Adderley - "Sack O' Woe" (1960)}
This soul jazz piece recorded live exemplifies the fusion of bebop virtuosity with Black sanctified church music through gospel call-and-response patterns. Adderley's alto saxophone timbre is "bright but thick," contrasting with Charlie Parker's grainier bebop tone while maintaining Parker's technical vocabulary. The 12-bar blues form features audience participation with yelling and clapping on beat 4, creating a preacher/congregation dynamic, while the piano takes on churchy gospel qualities. The rhythm shifts between swing feel and straight eights, emphasizing the groove and communal accessibility over bebop's pure intellectualism. This live recording demonstrates hard bop's goal of making jazz more communal and danceable while keeping bebop's sophisticated improvisational language.

\subsection*{Charles Mingus - "Better Git It in Your Soul" (1959)}
This hard bop composition exemplifies "organized chaos" through structured collective improvisation influenced by both bebop complexity and Duke Ellington's orchestral approach. The gospel-influenced piece features multiple horns creating dense polyphonic textures with a driving triplet feel and prominent trombone stabs. Mingus generates a big orchestral sound from a small group through careful compositional structure that allows simultaneous improvisation without losing coherence. The sanctified church quality and gospel influences demonstrate hard bop's reclamation of Black cultural roots after cool jazz. This piece extends bebop's virtuosity into a more compositionally sophisticated approach that balances structure with improvisational freedom.

\subsection*{John Coltrane - "Giant Steps" (1959)}
This piece represents the absolute pinnacle and endpoint of bebop's chord-based harmonic complexity through the revolutionary "Coltrane changes." The harmonic structure divides the octave into three equal parts with rapid ii-V chord progressions linking them, played at breakneck tempo. Coltrane's tenor saxophone features his "sheets of sound" technique with rapid cascading note patterns that outline every chord tone. Pianist Tommy Flanagan notably struggles during his solo because he hadn't rehearsed these difficult changes, demonstrating the extreme technical demands. After recording this ultimate extension of bebop harmony, Coltrane immediately moved toward modal and free jazz, making this the historical turning point where bebop's harmonic approach reached its limit.

\section*{Key Terms}

\textbf{Soul Jazz} - Funky, gospel-influenced hard bop substyle (Silver) \\
\textbf{Sheets of Sound} - Coltrane's rapid note cascades \\
\textbf{Walking Bass} - Steady quarter notes outlining chords \\
\textbf{Press Roll} - Blakey's signature drum fill between solos \\
\textbf{Organized Chaos} - Mingus's structured collective improvisation \\
\textbf{Contrafact} - New melody over existing chord changes ("Oleo" = "I Got Rhythm")

\section*{Readings}

\textbf{Monson, \emph{Saying Something} Ch. 3} - Jazz improvisation as conversation, rhythm section communication, call-and-response in African American music \\
\textbf{Szwed, \emph{Jazz 101} Ch. 20} - Hard bop as "subconsciously racial," mid-50s Black urban music from folk/church/blues

\end{document}
